\section{Existence for every index}\label{sec:existence}
The local lifting method used for Champernowne prefixes in the repository's
A075002 manuscript~\cite{prefix} also applies to a prefix starting at $n+1$.
The essential extra observation is that when $3\mid n$, the removed
triangular sum $n(n+1)/2$ vanishes modulo~3. We give the full argument
so that no certificate from another sequence is needed.

\begin{theorem}[All-index existence]\label{thm:existence}
For every positive integer $n$, there are infinitely many $t>n$ such that
\[
 n\mid F_n(t),\qquad F_n(t)=(n+1)\Vert(n+2)\Vert\cdots\Vert t.
\]
In particular, $a(n)$ exists for every $n\ge1$.
\end{theorem}

First put $F_n(n)=0$. In a fixed $d$-digit interval, let $B=10^d$ and
suppose $t+1,\ldots,t+L$ all lie in that interval. Appending these blocks
gives, for $L\ge0$,
\begin{equation}\label{eq:shift}
 F_n(t+L)-F_n(t)=(B^L-1)F_n(t)+tS_L+T_L,
\end{equation}
where $S_L=\sum_{j=0}^{L-1}B^j$ and
$T_L=\sum_{j=1}^L jB^{L-j}$. The telescoping identity
\begin{equation}\label{eq:weighted}
 (B-1)T_L=BS_L-L
\end{equation}
will be used only when the indicated division is legitimate.

\begin{lemma}[Two local shifts]\label{lem:shifts}
Assume the fixed-width condition in \eqref{eq:shift}.
\begin{enumerate}
\item Let $p\ge7$ be prime with $p\nmid B(B-1)$, and let
$h=\operatorname{ord}_p(B)$. If $L=chp^a$, $a\ge0$, and $p\nmid c$, then
\[
 F_n(t+L)-F_n(t)\equiv-\frac{chp^a}{B-1}\pmod{p^{a+1}}.
\]
For $e\ge1$, every multiple of
$p^e\operatorname{ord}_{p^e}(B)$ preserves both $F_n(t)$ and $t$
modulo $p^e$, provided the appended blocks stay in the interval.
\item If $a=v_3(L)\ge1$, then
\[
 F_n(t+L)-F_n(t)\equiv Lt+\frac{L(L+1)}2\pmod{3^{a+1}}.
\]
Every multiple of $2\cdot3^e$ preserves both $F_n(t)$ and $t$
modulo $3^e$, under the same interval condition.
\end{enumerate}
\end{lemma}
\begin{proof}
For the first assertion, the binomial theorem (or LTE) gives
$B^{chp^a}\equiv1\pmod{p^{a+1}}$, since $B^h\equiv1\pmod p$.
Thus $S_L\equiv0$ because $B-1$ is a unit. Equations
\eqref{eq:shift}--\eqref{eq:weighted} give the formula.
For a multiple of $p^e\operatorname{ord}_{p^e}(B)$, the same identities
give $B^L=1$, $S_L=T_L=0$ modulo $p^e$ and $p^e\mid L$.

For the second assertion write $B=1+c$ with $9\mid c$. Expansion gives
\[
 S_L-L=\sum_{j\ge1}c^j\binom{L}{j+1},\qquad
 T_L-\frac{L(L+1)}2=\sum_{j\ge1}c^j\binom{L+1}{j+2}.
\]
The valuations of the summands are respectively at least
$2j+a-v_3(j+1)$ and $2j+a-v_3((j+1)(j+2))$, both at least $a+1$.
Here $v_3(j+1)\le j$, and
$v_3((j+1)(j+2))\le2j-1$ for $j\ge1$ (check $j=1$ directly;
for $j\ge2$, $3^{2j}>(j+1)(j+2)$).
Also $B^L\equiv1\pmod{3^{a+2}}$ by LTE. Substitute in
\eqref{eq:shift}. If $2\cdot3^e\mid L$, the displayed shift is
zero modulo $3^e$, and $t$ is preserved as well.
\end{proof}

\begin{proof}[Proof of Theorem~\ref{thm:existence}]
The case $n=1$ is immediate. Write
\[
 n=R\,3^e\prod_{i=1}^s p_i^{e_i},\qquad
 R=2^\alpha5^\beta,\qquad 7\le p_1<\cdots<p_s.
\]
Let $D=\operatorname{lcm}_i\operatorname{ord}_{p_i^{e_i}}(10)$,
with $D=1$ for an empty product. Choose $d\equiv1\pmod D$ sufficiently
large, so $d\ge\max(\alpha,\beta)$ and $10^{d-1}>n$.
Then $B=10^d$ is zero modulo $R$ and is congruent to 10 modulo
each $p_i^{e_i}$. All step sizes below are consequently independent of $d$.

Define $H_0=2\cdot3^e$ if $e>0$, and $H_0=1$ otherwise; recursively put
\[
 H_i=\operatorname{lcm}\bigl(H_{i-1},
       p_i^{e_i}\operatorname{ord}_{p_i^{e_i}}(10)\bigr).
\]
For $1\le a<e$ let $L_a=\operatorname{lcm}(2\cdot3^a,R)$.
For $0\le a<e_i$ put
\[
 J_{i,a}=\operatorname{lcm}
   \bigl(\operatorname{ord}_{p_i}(10)p_i^a,H_{i-1},R\bigr),
 \qquad M_0=R3^{\min(e,1)}.
\]
An explicit displacement bound is
\begin{equation}\label{eq:capacity}
 W=2\sum_{a=1}^{e-1}L_a+
       \sum_{i=1}^s(p_i-1)\sum_{a=0}^{e_i-1}J_{i,a},
 \qquad 9\cdot10^{d-1}>M_0+W.
\end{equation}
Empty sums are zero. Increase $d$ within its progression to ensure
\eqref{eq:capacity}, and let $t_0$ be the least multiple of $M_0$
at least $10^{d-1}$. Every candidate below lies between $t_0$ and
$t_0+W<10^d$, which proves every required fixed-width condition.

The last block gives $F_n(t_0)\equiv t_0=0\pmod R$.
If $e>0$, then $3\mid n,t_0$, so
\[
 F_n(t_0)\equiv\sum_{j=n+1}^{t_0}j
   =\frac{t_0(t_0+1)-n(n+1)}2\equiv0\pmod3.
\]
At a root modulo $3^a$, $1\le a<e$, choose among endpoints
$t+jL_a$, $j=0,1,2$. Write $L_a=2c3^a$, where $3\nmid c$.
Lemma~\ref{lem:shifts} gives
\[
 \frac{F_n(t+jL_a)-F_n(t)}{3^a}
       \equiv jc(2t+1)\pmod3.
\]
Since $3\mid F_n(t)$ and $3\mid n$, the triangular identity implies
$t\equiv0$ or $2\pmod3$. The slope $c(2t+1)$ is therefore a unit.
Exactly one choice raises the root to $3^{a+1}$. All endpoints remain
multiples of $R$, preserving the last-block condition. This installs
$R3^e$; the period $H_0$ preserves the completed 3-part.

Now process $p_i$ in increasing order. Every prime divisor of
$p_j^{e_j}\operatorname{ord}_{p_j^{e_j}}(10)$ is at most $p_j$,
since the order divides $p_j^{e_j-1}(p_j-1)$. Thus $p_i\nmid RH_{i-1}$.
Consequently $J_{i,a}=chp_i^a$ with $h=\operatorname{ord}_{p_i}(10)$
and $p_i\nmid c$. At a root modulo $p_i^a$ (vacuous for $a=0$),
Lemma~\ref{lem:shifts} shows that the endpoints
$t+jJ_{i,a}$, $0\le j<p_i$, sweep all next $p_i$-adic digits,
with unit normalized slope $-ch/(B-1)$ modulo $p_i$.
Choose the unique root modulo $p_i^{a+1}$. The multiples of $H_{i-1}$
preserve all earlier odd prime powers, and those of $R$ preserve
the last-block condition. After installing $p_i^{e_i}$, the period
$H_i$ preserves it too. Induction yields $n\mid F_n(t)$.

We have $n<t<10^d$, hence the explicit bound $a(n)\le t-n<10^d-n$.
There are arbitrarily large admissible $d$; their disjoint digit
intervals give infinitely many endpoints. This proves the theorem.
\end{proof}

\begin{corollary}[A uniform bound]\label{cor:uniform}
For every $n\ge2$,
\[
 1\le a(n)<10^n-n.
\]
\end{corollary}
\begin{proof}
Retain the notation in the proof, and put $K=\operatorname{lcm}(R,H_s)$.
Every $L_a$ and $J_{i,a}$ divides $K$, and
\[
 K\le 2n\prod_i\varphi(p_i^{e_i})\le2n^2.
\]
The elementary inequalities $p^f\ge1+f(p-1)$ and
$\prod_j(1+x_j)\ge1+\sum_j x_j$ for nonnegative $x_j$ give
\[
 2\max(e-1,0)+\sum_i(p_i-1)e_i\le n-1.
\]
Thus $W\le2n^2(n-1)$ and $M_0\le n$.
Let $b=\lceil\log_2 n\rceil$ and choose
\[
 d=1+D\lceil b/D\rceil.
\]
Then $d\ge b+1$, so $10^{d-1}\ge10^b\ge n^3>n$ and
$9\cdot10^{d-1}>n+2n^2(n-1)$. All conditions in the construction hold.
Also $D\le\prod_i\varphi(p_i^{e_i})\le n-1$ (including the
empty-product case). If $D\ge b$, then $d=D+1\le n$.
If $D<b$, then $d\le b+D\le2b-1\le n$; the last inequality holds
for $n\ge3$, and $n=2$ falls in the first case.
The constructed endpoint therefore satisfies $n<t<10^d\le10^n$.
\end{proof}

\begin{proposition}[Exact minimum for every index]\label{prop:algorithm}
Every value $a(n)$ can be determined in at most $n^3$ modular
single-block transitions, together with at most $n$ binary block jumps.
\end{proposition}
\begin{proof}
The case $n=1$ is immediate. For $n\ge2$, process digit intervals
in increasing order, starting with the integer $n+1$ and stopping before
$10^n$. Within a $d$-digit interval, the state
$(F_n(t)\bmod n,t\bmod n)$ lies in a set of size $n^2$ and evolves by
\[
 (x,y)\longmapsto(10^d x+y+1,\ y+1)\pmod n.
\]
Test the residue after each actual appended block. If it is zero,
the first such endpoint gives the exact minimum. If a state repeats
before a zero occurs, determinism implies that the remaining orbit
in this interval has no zero. Jump over the unused part of the interval
by binary exponentiation of its affine transition and proceed to the
next digit width. A repeated state occurs within $n^2$ transitions,
irrespective of whether $10^d$ is invertible modulo $n$.
There are at most $n$ digit intervals, and Corollary~\ref{cor:uniform}
guarantees a zero before they end. No guessed numerical search limit
or assumed general minimum formula is used.
\end{proof}

